% Owner: classification_e.  The proof does not import CDP2020 Theorem2.2.
\subsubsection*{Removing the negative-slope horizontal fixed part}

\begin{proposition}[E.2.9 (Proved): no horizontal divisor when all fibres are irreducible]
Let $X$ be a compact complex threefold diffeomorphic to $S^6$,
and let $f:X\to\mathbb P^1$ be a surjective holomorphic map
with connected fibres.  If every reduced fibre is irreducible,
then $X$ has no horizontal prime divisor for $f$.
\end{proposition}
\begin{proof}
Suppose that $H$ is such a divisor.  It is Cartier because
$X$ is smooth.  Since $H$ is horizontal and every reduced
fibre $S$ is irreducible, $H$ meets $S$ and does not contain
it.  As in E.2.8, restriction of its canonical section is a
non-zero-divisor on every local branch of $S$.

We first record the elementary vanishing used below.  On a
compact reduced irreducible complex space $S$, let $D$ be a
nonzero effective Cartier divisor and $V$ a vector bundle.
Multiplication by the powers of the canonical section
identifies
\[
 H^0(S,V(-nD))\subset H^0(S,V),\qquad n\ge0,
\]
with a decreasing sequence of finite-dimensional vector
spaces.  A section in their intersection vanishes to every
order along $D$.  Krull intersection in a local ring at a
point of $D$ makes its germ zero.  After pullback to the
irreducible normalization, analytic continuation makes the
section zero everywhere.  The decreasing sequence therefore
stabilizes at zero.  In particular
\[
 H^0(S,V(-nD))=0\quad\hbox{for all sufficiently large }n.
\]
The same conclusion holds for a vector bundle on a thickened
fibre $mS$.  Indeed $S$ is Cartier in $X$, and the finite
filtration by powers of its ideal has successive quotients
\[
 V|_S\otimes\mathcal O_S(-jS),\qquad 0\le j<m.
\]
Apply the preceding reduced-space statement to these finitely
many vector bundles, with the same nonzero divisor $H|_S$,
and then use their exact sequences.  This also proves that
the multiplication maps used here remain injective on $mS$:
its only associated component is $S$ and $H$ contains none
of its local branches.

Let $B$ be any holomorphic vector bundle on $X$.  The map
$f$ is flat, since its local rings are torsion-free over the
discrete valuation rings of the smooth base.  Every analytic
fibre $F_t$ is a Cartier, Cohen--Macaulay, Gorenstein surface,
possibly nonreduced.  Its dualizing line is
\[
 \omega_{F_t}=K_X|_{F_t},
\]
where the constant one-dimensional factor from the target
has been trivialized.  Serre duality on this compact
Gorenstein surface gives
\[
 H^2(F_t,B(nH)|_{F_t})^*
 \simeq H^0\bigl(F_t,(B^*\otimes K_X)(-nH)|_{F_t}\bigr).
\]
The compact-space duality being used includes singular
Cohen--Macaulay spaces; it is the absolute duality theorem
of \citet[\S3, Theorem~1, p.~81, and the complete-intersection
calculation on p.~91]{RR70push}, already specified in E.3.1,
not a smooth-fibre replacement.
The preceding filtration argument makes the right hand
side zero for all sufficiently large $n$, for each $t$.

For fixed $n$, the nonvanishing locus
\[
 Z_n=\{t\in\mathbb P^1:
        h^2(F_t,B(nH)|_{F_t})>0\}
\]
is closed by semicontinuity.  Multiplication on the dual
section spaces gives $Z_{n+1}\subset Z_n$.  Their intersection
is empty by fibrewise vanishing.  Compactness of $\mathbb P^1$
therefore supplies a single $N$ with $Z_N=\varnothing$.
Constant-zero cohomology and analytic cohomology-and-base-change
\cite[Chapter~III, Theorem~4.12(ii)]{BS76push} give
\[
 R^2f_*(B(NH))=0.
\]
Higher direct images vanish in degrees greater than two by
the fibre dimension.  The coherent sheaf $R^1f_*(B(NH))$
has vanishing first cohomology after a sufficiently positive
twist on $\mathbb P^1$.  The projection formula and Leray
sequence consequently give, for some integer $a\gg0$,
\[
 H^2(X,B\otimes M_+)=H^3(X,B\otimes M_+)=0,
 \qquad
 M_+=\mathcal O_X(NH)\otimes f^*\mathcal O(a).
\]
All coherence, semicontinuity and constant-zero base-change
uses are the proper analytic theorems audited in E.3.1;
the argument does not infer base change from coherence alone.

Independently, the reduced-space vanishing argument applied
to the connected smooth manifold $X$ and its nonzero divisor
$H$ gives
\[
 H^0(X,B\otimes M_-)=0,
 \qquad M_-=\mathcal O_X(-bH),\quad b\gg0.
\]
We must arrange these different vanishings for one line
bundle.  The exponential sequence gives
\[
 \operatorname{Pic}(X)\simeq H^1(X,\mathcal O_X)
\]
because $H^1(X,\mathbb Z)=H^2(X,\mathbb Z)=0$.
Let $\ell_-,\ell_+$ be the corresponding additive classes.
The affine line
\[
 \ell(z)=(1-z)\ell_-+z\ell_+,\qquad z\in\mathbb C,
\]
is realized by a holomorphic family of line bundles on
$X\times\mathbb C$: choose additive holomorphic Cech
cocycles representing the two classes and exponentiate their
linear combination.  Equivalently, use the integrable
Dolbeault operators obtained by a holomorphic linear
combination of fixed closed $(0,1)$-form representatives.
Thus no representability or universal-Picard-space assertion
is needed here.

For the proper projection $X\times\mathbb C\to\mathbb C$,
the cohomology jumping loci
\[
 T_j=\{z:h^j(X,B\otimes M_z)>0\},\qquad j=0,2,3,
\]
are closed analytic subsets of $\mathbb C$ by the proper
analytic finite-complex theorem
\cite[Chapter~III, Theorem~4.1]{BS76push}: the condition is
expressed by the rank minors in a local finite free complex.
In particular no closedness claim from the corrected
Grauert Section~7, Satz~2 is used.  The set $T_0$ omits $z=0$,
while $T_2$ and $T_3$ omit $z=1$.  Each is therefore proper,
and their finite union cannot fill the irreducible complex
line.  Choose $z$ outside that union.  Then
\[
 \chi(X,B\otimes M_z)=-h^1(X,B\otimes M_z)\le0.
\]
Every line bundle on $X$ is topologically trivial, so
Hirzebruch--Riemann--Roch makes this Euler characteristic
equal to $\chi(X,B)$.

Finally take $B=T_X$.  Since $H^2(X,\mathbb Z)=H^4(X,\mathbb Z)=0$,
the only contributing term in its Riemann--Roch formula is
\[
 \chi(X,T_X)=\frac12\langle c_3(X),[X]\rangle
            =\frac12\chi_{\mathrm{top}}(S^6)=1.
\]
This contradicts the preceding inequality.  Hence no
horizontal prime $H$ exists.
\end{proof}

\begin{proposition}[E.2.10 (Proved): persistence for either canonical slope]
Under the hypotheses of E.2.5, for every nonzero
$D=kl-3k-4l$ one has $E_0=0$ and $e=1$.
After shrinking the disc, the evaluation of
$\pi_*\mathcal L$ gives a proper flat holomorphic map
\[
 g:\mathcal X\longrightarrow\mathbb P^1\times\Delta
\]
whose restriction to every fibre is its intrinsic algebraic
reduction.  A smooth general central fibre of $g$ is a
complex two-torus.  Its normalized punctured period has
an interior limit as in E.2.7.

Two families satisfying E.2.5 cannot have the same smooth
central manifold unless their ordered multiplicity pairs
agree, regardless of the signs of their canonical slopes.
Uniqueness of a smooth completion at an interior period and
the arbitrary deformation-chain statement E.2.4 are not
asserted.
\end{proposition}
\begin{proof}
The case $D>0$ is E.2.6.  If $D<0$ and $E_0$ had a
horizontal prime component, the final part of E.2.8 would
give $b_1(F)=1$ for a smooth general fibre of $g_0$ and
would make every reduced fibre of $g_0$ irreducible.
E.2.9 would then forbid precisely that horizontal component.
Thus for either sign $E_0$ is vertical.

The connected-divisor-limit argument in the middle of the
proof of E.2.6 only uses verticality, not the sign of $D$.
Its generic choice of pencil section has central zero
support equal to the disjoint union of the vertical fixed
part and $e$ distinct regular fibres.  The total divisor
is reduced and irreducible with connected central support
after normalization and Stein factorization.  Consequently
$E_0=0$ and $e=1$.  Surjectivity of evaluation, flatness and
the torus assertion follow from the remainder of that same
proof.  The hypothesis of E.2.7 is now satisfied for either
sign, and gives the interior period limit.

Suppose two such one-sided families have the same central
manifold $X_0$.  Its algebraic reduction is intrinsic, so
both pencil lines are the same line
$L=g_0^*\mathcal O(1)$ up to isomorphism.  With indices
$i=1,2$, the two root identities give
\[
 K_{X_0}^{A_i}=L^{D_i},\qquad
 L^{D_1A_2-D_2A_1}=\mathcal O.
\]
The line $L$ is not torsion, because a positive power has
a nonzero section with zeros.  Hence $D_1/A_1=D_2/A_2$.
E.2.3 recovers the ordered coprime multiplicity pair from
this rational number.  This proves the common-limit
assertion without identifying the limit with a member of
the given period construction.
\end{proof}

The vanishing argument is an application of standard proper
cohomology, duality and Riemann--Roch machinery; no novelty
claim is intended.  It proves the needed absence of a
horizontal fixed divisor under the rank-one hypothesis.
It does not establish the different claim that every
holomorphic Hopf-surface fibration on a complex six-sphere
is impossible.
